\documentclass[11pt,a4paper]{article}
\usepackage[utf8]{inputenc}
\usepackage[T1]{fontenc}
\usepackage[english]{babel}
\usepackage{amsmath,amssymb,amsthm}
\usepackage{mathtools}
\usepackage[colorlinks=true,linkcolor=blue,citecolor=blue,urlcolor=blue]{hyperref}
\usepackage[numbers]{natbib}
\usepackage[most]{tcolorbox}
\usepackage{geometry}
\geometry{margin=2.5cm}
\usepackage{booktabs}
\usepackage{longtable}
\usepackage{array}

\pdfstringdefDisableCommands{%
\let\ensuremath\relax
}

\newtcolorbox{argumentbox}[1]{colback=gray!5,colframe=black!60,fonttitle=\bfseries,title=#1}
\newtcolorbox{resolutionbox}[1][]{colback=green!5,colframe=green!40!black,fonttitle=\bfseries,title=Resolution #1}
\newtcolorbox{openprobbox}[1][]{colback=red!5,colframe=red!60!black,fonttitle=\bfseries,title=Open Problem #1}

\title{\textbf{On the Necessity of the Binary Distinction and the Dissolution of the Bearer Problem}\texorpdfstring{\\[4pt]\large An Addendum to \emph{Existence as Pulsating Closure} (D19)}{: An Addendum to Existence as Pulsating Closure (D19)}\\[6pt]
{\large PDL Document D19ad}}

\author{C\'edric Laubscher\\[4pt]
\small Independent Researcher, Switzerland\\
\small \href{https://cedriclaubscher.ch}{cedriclaubscher.ch}}

\date{July 2026}

\begin{document}
\maketitle

\begin{abstract}
D19~\citep{D19} establishes that a reproducible distinction is the minimal tenable candidate for existence, and that instantaneous or featureless configurations are logically indistinguishable from nothingness. What D19 states but does not derive is the specific arithmetic content of axiom C1: why the minimal distinction must be \emph{binary} rather than admitting three or more states, and what must already be presupposed for a trace to constitute a memory rather than an isolated event. This addendum closes both gaps. First, we show that the cardinality of the minimal distinction is not a free choice but a forced consequence of the nature of a single act of distinguishing, considered independently of any prior commitment to plurality. Second, we identify a residual presupposition in the trace-to-memory argument of D19 --- the existence of a bearer that persists between inscription and consultation --- and show that this presupposition dissolves once the reality of the discrete instant itself is denied, rather than requiring an additional axiom. The argument proceeds entirely within the philosophical register of D19 and DN~\citep{DN}; no new physical or combinatorial result is claimed.
\end{abstract}

\clearpage

\section{Motivation and Scope}
\label{sec:motivation}

D19 argues that existence, reduced to its logical minimum, is nothing other than a reproducible distinction: an alternation between two states capable of returning to itself~\citep[\S 2--3]{D19}. The same reduction is repeated, in narrative form, in Chapter 1 of DN (\emph{Nothing, then Something})~\citep{DN}, with an additional argument --- absent from D19 --- for why the alternation must be \emph{discrete} rather than continuous: a continuous process admits no guarantee of exact return, whereas a cycle that must close without drift is, by that very requirement, discrete (DN, Chapter 1, \S \emph{A Granular Universe by Necessity}).

Two steps in this reduction are asserted rather than derived in both texts. This addendum supplies the missing derivations. Neither step introduces a new axiom, a new combinatorial object, or a new physical claim: both remain strictly upstream of C1, in the same register as D19 itself.

\section{Why Two: The Boundary as a Single Act}
\label{sec:why_two}

\subsection{The Gap in D19 and DN}

D19 introduces the binary alternation directly: ``an elementary alternation between two mutually exclusive states''~\citep[\S 2]{D19}. DN does the same: ``imagine two possible states: agreement and disagreement, yes and no, $+1$ and $-1$'' (DN, Chapter 1, \S \emph{The Minimal Pulse}). In neither text is the number two itself derived from anything more primitive. A reader is entitled to ask: why not three states, or a continuum of states subsequently rendered discrete for other reasons?

\subsection{The Argument}

\begin{argumentbox}{Step 1 --- Distinction as a Single Undecomposed Act}
Consider the most minimal act compatible with the emergence of any difference at all: not the positing of several pre-existing terms placed in relation, but a single, undecomposed act of distinguishing --- the drawing of one boundary, prior to any content on either side of it~\citep{SpencerBrown1969}. This is a strictly weaker assumption than positing ``two things side by side'': it does not presuppose plurality, only the possibility of a difference being marked at all.
\end{argumentbox}

\begin{argumentbox}{Step 2 --- What a Boundary Is Independently of Content}
A boundary, considered purely as a boundary and not yet as separating any particular content, has exactly two sides by definition: an inside and an outside, a marked and an unmarked state. This is not a further stipulation added to the act of distinguishing; it is what the word \emph{boundary} means. A single act of distinguishing that produced three or more sides would not be a single boundary; it would already be a compound structure requiring several acts, and hence would not be the minimal case under consideration.
\end{argumentbox}

\begin{resolutionbox}
The cardinality two is therefore not chosen among several admissible options; it is the unique cardinality compatible with there being \emph{one} act of distinction rather than several. Three states would require either a second, independent boundary (already a plural, non-minimal construction) or a single boundary with more than two sides, which is not a coherent notion of boundary at all. The binary character of C1's alternation is accordingly forced by the requirement of minimality itself, not imposed as an additional constraint alongside it.
\end{resolutionbox}

\subsection{Relation to DN's Discreteness Argument}

DN's argument for discreteness (\S\ref{sec:motivation} above) is independent of, and compatible with, the argument just given. DN shows that \emph{given} some alternation of states, that alternation must be discrete for a cycle to close exactly. The present section shows, prior to that, why the alternation is binary rather than of higher cardinality. The two arguments are complementary and jointly close the derivation of C1's `two discrete states' from the bare requirement of a single, minimal, reproducible distinction.

\section{The Bearer Problem and Its Dissolution}
\label{sec:bearer}

\subsection{The Gap: What a Trace Presupposes}

Both D19 and DN argue that an event with no possibility of recurrence is indistinguishable from non-occurrence, and hence that existence requires repeatability~\citep[\S 2]{D19}(DN, Chapter 1, \S \emph{Instantaneous Existence Does Not Count}). This argument, examined closely, presupposes more than it states. A mark or trace left by an event becomes a \emph{memory} --- rather than a brute fact indistinguishable from any other --- only if it is read, at a later point, \emph{as referring to} an earlier point. This requires two things to be true simultaneously at the moment of reading: the trace itself, and something that persists between the moment of inscription and the moment of consultation, such that the consultation is a reading of \emph{that} trace rather than the accidental encounter of two unrelated facts.

Call this persisting requirement the \emph{bearer}. Neither D19 nor DN names or examines this presupposition; the repeatability argument is stated as though it were self-sufficient, when in fact it silently imports a notion of persistent identity across two moments --- precisely the kind of temporal structure the argument is meant to explain rather than assume.

\subsection{Why the Bearer Cannot Simply Be Postulated}

One might attempt to discharge this presupposition by treating the bearer as an additional primitive: some minimal substrate that persists while its states change. This move should be resisted. Positing a persisting substrate as a primitive is exactly the kind of unexplained carry-over from ordinary metaphysics that D19 and DN set out to avoid at the outset (D19, \S 1: ``no space, no time, no matter, and no antecedently specified inventory of entities''); it would smuggle back in, under a new name, the very continuant ontology the reduction to pulsating closure is designed to do without.

\subsection{The Dissolution}

\begin{argumentbox}{The Move}
The apparent need for a bearer arises only if one presupposes two separated instants --- an instant of inscription and a later, distinct instant of consultation --- with an interval between them across which something would have to persist unchanged in order to bridge the gap. This presupposition of separated instants is precisely what a reproducible distinction, prior to any further structure, does not yet contain. There is, at this stage, no instant in the strong sense: no punctual, self-standing ``now'' detachable from the rest and capable of being placed in a before/after relation to another such ``now''.
\end{argumentbox}

\begin{resolutionbox}
Once the reality of the discrete instant is denied at this level, the bearer problem does not receive an answer; it is dissolved. There is no gap between inscription and consultation for something to bridge, because the picture of two separated instants was never available in the first place. What remains is simply the continuous, undivided ``something'' with which the reduction begins in D19 \S 2 and DN Chapter 1 --- not a new entity introduced to solve the bearer problem, but the same primitive already in play, now recognised as sufficient without supplementation. The bearer is not found; the question that seemed to require one is shown to rest on a presupposition the framework never needed to grant.
\end{resolutionbox}

\subsection{A New Question Opened, Not Closed}

This dissolution is not free of cost. If no discrete instant is available prior to the first distinction, a further question arises: how does the first distinction --- which, on the reading of Section~\ref{sec:why_two}, is a single act --- itself occur, given that an act appears, on the face of it, to require a before and an after in order to be an act at all? This tension is real and is not resolved here.

\begin{openprobbox}
\label{op:act_without_instant}
Characterise, within the PDL framework, how a single act of distinction can be said to occur in the absence of any prior discrete instant, without presupposing the very temporal structure (before/after, an instant of occurrence) that the act is supposed to inaugurate.
\end{openprobbox}

This question is philosophical in the strict sense: it does not reduce to a combinatorial or numerical claim about signed graphs, and no attempt is made here to force it into that register. It is recorded as an open problem of the same family as those D19 leaves explicitly unresolved (D19, \S 4, on whether $\varepsilon$ carries genuine transcendence or is a bookkeeping device)~\citep{D19}.

\section{Summary of the Two Closures}
\label{sec:summary}

\begin{table}[h]
\centering
\caption{The two gaps addressed by this addendum, and their resolution.}
\label{tab:summary}
\begin{tabular}{p{3.6cm}p{5.2cm}p{5.6cm}}
\toprule
& \textbf{D19 / DN (as published)} & \textbf{This addendum} \\
\midrule
Why two states & Asserted directly & Derived: forced by the nature of a single, undecomposed act of distinction (\S\ref{sec:why_two}) \\
Why discrete & Derived (DN only): exact closure requires it & Unaffected; retained as complementary \\
Why repeatable & Derived: non-recurrence is indistinguishable from non-occurrence & Unaffected; retained \\
The bearer & Presupposed silently & Named, and dissolved by denying the discrete instant prior to the first distinction (\S\ref{sec:bearer}) \\
\bottomrule
\end{tabular}
\end{table}

\section{Status and Relation to the Rest of the Corpus}
\label{sec:status}

Both arguments of this addendum are philosophical derivations in the register of D19, not combinatorial or physical results. They neither strengthen nor weaken any numerical claim of the corpus's mode-A programme (D01--D45 and following); C1 is not thereby proved in any stronger technical sense than before. What changes is the epistemic status of C1's origin: the two-state, discrete, repeatable character of the minimal pulsation is shown to follow from a single, weaker requirement --- one undecomposed act of distinction, without an available prior instant --- rather than being a felicitous stipulation later found to generate a rich physics. This is relevant, in particular, to any future reading of the PDL-V programme (DL01, DL02) and of the language-oriented exploration reported separately (DL03, forthcoming): both inherit whatever robustness C1 itself carries, and this addendum is offered as a contribution to that robustness, not as a new axiom or a new layer of the technical corpus.

Whether Open Problem~\ref{op:act_without_instant} admits a resolution within the philosophical register alone, or requires a technical construction (for instance, a formal treatment of ``act'' that does not presuppose a prior instant), is left unaddressed here.

\clearpage

\bibliographystyle{unsrt}
\bibliography{D19ad_references}

\end{document}
